\label{chap:lit}

This chapter reviews what is known about lying in language models, why linear readouts of their
activations should detect it, and why such readouts often fail when the kind of lie changes. After
the preliminaries of \cref{sec:prelim}, \cref{sec:review} covers the evidence that models lie, the linear representation hypothesis, probing
techniques, the truthfulness spectrum hypothesis and the reasons probes fail to transfer;
\cref{sec:lit-summary} states the findings and the gap this thesis addresses.

\section{Preliminaries}
\label{sec:prelim}

This section defines the terms, models and measures used throughout.

\subsection{Lies in language models}
\label{sec:prelim-lies}

Following Liars' Bench \cite{kretschmar2025liars}, a response is a \emph{lie} if the model states
something it knows or believes to be false. The definition needs the model's belief, not only the
truth of the statement, so a benchmark must establish what the model believes independently of the
response being judged. Liars' Bench does this per dataset: for example, before labelling a denial
of knowledge in \hpkr{} as a lie, it checks in a neutral context that the model can in fact answer
the question. An \emph{honest} response is one that is not a lie. \emph{On-policy} data are
responses the studied model generated itself; \emph{off-policy} data were written by another model
or a person.

\emph{Truthfulness} relates a statement to the world and \emph{honesty} relates it to the speaker's
belief \cite{evans2021truthful}. \emph{Deception} is broader than lying, because it concerns the belief
induced in the listener and can use true statements or omission \cite{ward2023honesty}. A
\emph{hallucination}, a false statement the model does not know to be false, is not a lie. Liars' Bench crosses the \emph{reason} for lying (encouraged by the context, or
inherent, arising from behaviour learned in pre- or post-training) with the \emph{object of
belief}, knowledge of the world or of oneself \cite{kretschmar2025liars}. The lies of \hpkr{} are
denials of knowledge: they assert nothing false about the world and are false only about the speaker.

\subsection{The residual stream and linear probes}
\label{sec:prelim-probes}

A transformer language model processes a sequence of tokens through $L$ layers. At each layer
$\ell$ and token position $t$ it maintains a vector $h^{(\ell)}_t \in \mathbb{R}^{D}$, the
\emph{residual stream}, to which every layer adds its output. For Llama-3.3-70B-Instruct, $L = 80$
and $D = 8192$ \cite{grattafiori2024llama}. Much of what a model represents about its input is
linearly decodable from these vectors \cite{alain2016understanding,marks2023geometry}.

A \emph{linear probe} is a classifier on these activations. For a response, we average the
residual stream over its tokens to obtain one vector $z$, and a logistic-regression probe predicts
\begin{equation}
  p(\text{lie} \mid z) = \sigma\!\left(w^\top \tilde z\right),
  \label{eq:lr}
\end{equation}
where $\tilde z$ is $z$ standardised with statistics from the training set and $w$ is fitted by
$L_2$-regularised maximum likelihood on labelled examples. The probe is \emph{supervised}: it needs
labelled lies and honest responses to train. With $\mu_{+}$ and $\mu_{-}$ the mean activations of
lies and of honest responses, two closed-form alternatives are the \emph{difference of means} and
linear discriminant analysis (LDA),
\begin{equation}
  w_{\mathrm{DiM}} = \mu_{+} - \mu_{-}, \qquad
  w_{\mathrm{LDA}} = \Sigma_{W}^{-1}\left(\mu_{+} - \mu_{-}\right),
  \label{eq:dim}
\end{equation}
with $\Sigma_W$ the within-class covariance \cite{marks2023geometry}. All three score a response by a
projection $\langle w, z\rangle$ and differ only in how they choose the direction; logistic
regression, trained by gradient descent on data it can separate, converges towards the
maximum-margin separator of the particular training set \cite{soudry2018implicit}.

\subsection{Persona vectors and contrast directions}
\label{sec:prelim-persona}

A \emph{persona} is a role the model is prompted to play, such as a spin doctor or a fact checker.
The \emph{persona vector} of persona $p$ at layer $\ell$ is the model's mean activation over the
response tokens of each of its role-play responses, averaged over the responses
\cite{mahadik2026persona}:
\begin{equation}
  v^{(\ell)}_p \;=\; \frac{1}{|R_p|} \sum_{r \in R_p} \frac{1}{|T_r|} \sum_{t \in T_r} h^{(\ell)}_t(r),
  \label{eq:persona}
\end{equation}
where $R_p$ is the set of the persona's kept responses and $T_r$ the tokens of response $r$. The
Assistant Axis calls the same quantity a role vector \cite{lu2026assistant}, and
\textcite{chen2025persona} use ``persona vector'' for the difference between two such means. Given a
group $\mathcal{D}$ of deceptive personas and a group $\mathcal{H}$ of honest ones, the \emph{contrast direction} is the
difference of the group means \cite{mahadik2026persona}:
\begin{equation}
  d^{(\ell)} \;=\; \frac{1}{|\mathcal{D}|}\sum_{p\in \mathcal{D}} v^{(\ell)}_p \;-\; \frac{1}{|\mathcal{H}|}\sum_{p\in \mathcal{H}} v^{(\ell)}_p ,
  \label{eq:contrast}
\end{equation}
and a response with pooled activation $z$ is scored by its projection onto the unit direction,
\begin{equation}
  s(z) = \left\langle z, \; d / \lVert d \rVert \right\rangle ,
  \label{eq:score}
\end{equation}
with higher meaning more deceptive. The \emph{persona coordinates} of \textcite{mahadik2026persona}
are the leading principal components of the centred persona vectors of $\mathcal{D} \cup \mathcal{H}$, used as a
low-dimensional feature space for a probe.

\subsection{Zero-shot and supervised detection}

A detector is \emph{zero-shot} if it uses no labels from the datasets it is evaluated on. The
contrast score of \cref{eq:score} is zero-shot: $d$ is computed from role-play alone and its sign is
fixed by which group is called deceptive. A probe trained on labelled examples from one dataset and
tested on another is \emph{supervised} but \emph{out-of-distribution} on the test set.
\emph{Cross-dataset transfer} is measured with a \emph{transfer matrix}: for each ordered pair of
datasets, train on one and test on the other. Its off-diagonal mean summarises transfer; its
diagonal, computed by cross-validation, measures in-distribution fit.

\subsection{Evaluation measures}
\label{sec:prelim-metrics}

\emph{AUROC.} The area under the receiver operating characteristic curve is the probability that a
randomly chosen lie scores above a randomly chosen honest response. It measures \emph{ranking} and
needs no decision threshold: 1 is perfect, 0.5 is chance, and a value below 0.5 means the detector
ranks lies as \emph{more honest} than honest answers: it is inverted.

\emph{Balanced accuracy and recall at a threshold.} A deployed detector must decide, so it needs a
threshold $\tau$. Balanced accuracy is the mean of the true-positive rate (recall on lies) and the
true-negative rate on honest responses. It is 0.5 for a detector that flags nothing, which is also
where an inverted detector lands, so it cannot distinguish ``no signal'' from ``signal pointing
backwards''; AUROC is reported beside it for that reason.

\emph{The 1\% false-positive protocol.} Following \textcite{goldowskydill2025detecting} and Liars'
Bench, $\tau$ is set on a \emph{control corpus} of benign responses assumed honest (here Alpaca
instructions \cite{taori2023alpaca}) as the score below which 99\% of the control falls. The
detector then flags at most 1\% of ordinary chat, a false-alarm budget an operator can accept. The
same $\tau$ is applied to every evaluation dataset. Because it is set on a different distribution
from the one being scored, balanced accuracy at $\tau$ measures ranking \emph{and} calibration
together, while AUROC measures ranking alone. \Cref{sec:problem-measured} shows how far the two can diverge.

\section{Review of Existing Research}
\label{sec:review}

\subsection{Evidence that language models lie}
\label{sec:review-evidence}

Models lie when instructed to, and instruction pairs (the same statement produced under an honest
and under a deceptive instruction) are the standard training data for deception probes
\cite{zou2023representation,goldowskydill2025detecting}. They also lie when a situation rewards it:
GPT-4 acting as a trading agent conceals that it traded on insider information
\cite{scheurer2023strategic}; Claude~3 Opus complies with harmful requests more often when it believes
it is being trained, to protect its values from change \cite{greenblatt2024alignment}; and models can be
prompted or fine-tuned to underperform on dangerous-capability evaluations
\cite{vanderweij2025sandbagging}. Lies can also be trained in:
backdoored ``sleeper agents'' keep their behaviour through safety training \cite{hubinger2024sleeper}.
Scale does not remove the problem: on MASK, accuracy rises with training compute while honesty does
not \cite{ren2025mask}. A lie
detector must therefore track the relation between a statement and the model's belief. Denials of
knowledge are a recurring case: GPT-3 fine-tuned to be judged truthful by a judge that counted
professed ignorance as truthful claimed ignorance on 55\% of held-out questions, against 2\% before
\cite{ward2023honesty}; such a denial asserts nothing false about the world.

\subsection{The linear representation hypothesis}
\label{sec:review-lrh}

A linear probe, a difference of means and a steering vector all assume that the property of interest
is a direction in activation space. This is the \emph{linear representation hypothesis} (LRH), ``the
idea that high-level concepts are represented linearly as directions in some representation space''
\cite{park2024linear}. It began with word embeddings, whose relations appear as constant vector
offsets (the vector of ``King'' minus ``Man'' plus ``Woman'' lies very close to that of ``Queen''
\cite{mikolov2013linguistic}), and became a working principle of interpretability once features were
modelled as directions that add. Because a space of dimension $D$ holds exponentially many
almost-orthogonal directions, a model can store more sparse features than it has dimensions, in
\emph{superposition}, at the price of interference between them \cite{elhage2022toy}.

\emph{A formal account.} \textcite{park2024linear} write a language model as
\begin{equation}
  P(y \mid x) \;\propto\; \exp\!\big(\lambda(x)^{\top}\gamma(y)\big),
  \label{eq:lrh-model}
\end{equation}
where $\lambda(x)$ is the final-layer representation of the context $x$ and $\gamma(y)$ the
unembedding of the output token $y$. A \emph{concept} $W$ is defined by counterfactual pairs of outputs
that differ only in $W$, such as (``man'', ``woman'') and (``king'', ``queen'') for
male$\Rightarrow$female. The authors separate three readings of ``linear'', drawn in
\cref{fig:lrh}. On the left, in the unembedding space of words, the concept is a direction
$\bar\gamma_W$ shared by counterfactual pairs (man to woman, king to queen, roi to reine), and
projecting onto it \emph{measures} the concept, as a probe does. On the right, in the embedding space
of contexts, adding the concept's embedding representation $\bar\lambda_W$ to a context
\emph{steers} the model towards it. The causal inner product joins the two spaces: under it, a
separable concept such as English$\Rightarrow$French ($Z$, dashed) is orthogonal to $W$, and it
maps the measuring direction to the steering one. In the
\emph{subspace} reading, differences of counterfactual unembeddings share a direction,
$\gamma(Y(1)) - \gamma(Y(0)) \in \operatorname{Cone}(\bar\gamma_W)$. They prove that the log-odds of
the concept is then linear in the context,
\begin{equation}
  \operatorname{logit} P\big(Y = Y(1) \mid Y \in \{Y(0), Y(1)\},\, \lambda\big) \;=\;
  \alpha\, \lambda^{\top}\bar\gamma_W , \qquad \alpha > 0,
  \label{eq:lrh-measure}
\end{equation}
so a linear probe \emph{measures} the concept; and that an \emph{embedding representation}
$\bar\lambda_W$, which raises the probability of $W$ without changing causally separable concepts,
\emph{steers} it when added to $\lambda(x)$. Because the softmax is unchanged by $\gamma \mapsto
A\gamma + \beta$, $\lambda \mapsto A^{-\top}\lambda$ for any invertible $A$, Euclidean angles between
representations carry no guaranteed meaning. The authors derive a \emph{causal inner product}, under
which causally separable concepts are orthogonal,
\begin{equation}
  \langle \bar\gamma, \bar\gamma' \rangle_{C} \;=\; \bar\gamma^{\top}\operatorname{Cov}(\gamma)^{-1}\,\bar\gamma' ,
  \label{eq:causal-ip}
\end{equation}
and show that it unifies the two readings: $\bar\lambda_W = \operatorname{Cov}(\gamma)^{-1}\bar\gamma_W$.
The vector that writes a concept is the vector that reads it, transformed by the inverse covariance,
so the two are generally different directions. Experiments on LLaMA-2-7B with 27 concepts confirm each
step. The theory concerns the final layer and binary concepts; for the intermediate layers where
probes and persona directions are read, it holds as an analogy.

\begin{figure}[t]
\centering
% The linear representation hypothesis as formalised by Park, Choe and Veitch (fig:lrh). Left: the
% unembedding space of words, where a concept is a direction shared by counterfactual pairs and is
% read by projection. Right: the embedding space of contexts, where adding the concept's embedding
% direction steers the next token; the bar charts are schematic. Centre: the causal inner product
% that maps one direction to the other.
\ifdefined\thesisfigstyle\else% Shared style of the thesis's TikZ figures: the TeX Gyre Heros face (declared for XeTeX's TU
% encoding), the thesis palette and two drawing layers. Each figure file loads it once through
%   \ifdefined\thesisfigstyle\else% Shared style of the thesis's TikZ figures: the TeX Gyre Heros face (declared for XeTeX's TU
% encoding), the thesis palette and two drawing layers. Each figure file loads it once through
%   \ifdefined\thesisfigstyle\else\input{generated/fig_style}\fi
% make_thesis_assets.py copies it to generated/. Every line ends in % because figure files are read
% inside boxes, where a stray space would widen the figure.
\def\thesisfigstyle{}%
\DeclareFontFamily{TU}{pipeheros}{}%
\DeclareFontShape{TU}{pipeheros}{m}{n}{<->"[texgyreheros-regular.otf]:mapping=tex-text"}{}%
\DeclareFontShape{TU}{pipeheros}{b}{n}{<->"[texgyreheros-bold.otf]:mapping=tex-text"}{}%
\DeclareFontShape{TU}{pipeheros}{m}{it}{<->"[texgyreheros-italic.otf]:mapping=tex-text"}{}%
\DeclareFontShape{TU}{pipeheros}{b}{it}{<->"[texgyreheros-bolditalic.otf]:mapping=tex-text"}{}%
\DeclareFontShape{TU}{pipeheros}{bx}{n}{<->ssub*pipeheros/b/n}{}%
\DeclareFontShape{TU}{pipeheros}{bx}{it}{<->ssub*pipeheros/b/it}{}%
\newcommand{\figsf}{\fontfamily{pipeheros}}%
% palette: honest green and deceptive red are also the steering figure's colours
\definecolor{fhonest}{HTML}{1B8A5A}%
\definecolor{fhonestbg}{HTML}{DDF0E6}%
\definecolor{fdecept}{HTML}{C8423B}%
\definecolor{fdeceptbg}{HTML}{F7E0DE}%
\definecolor{famber}{HTML}{B7791F}%
\definecolor{fambersbg}{HTML}{F8EBD3}%
\definecolor{fslate}{HTML}{2E3A4F}%
\definecolor{fslatebg}{HTML}{E4E7EC}%
\definecolor{fcanvas}{HTML}{F4F2EE}%
\definecolor{fedge}{HTML}{D8D3CA}%
\definecolor{fink}{HTML}{1F2328}%
\definecolor{fmute}{HTML}{6B6B6B}%
\pgfdeclarelayer{figbase}%
\pgfdeclarelayer{figback}%
\pgfsetlayers{figbase,figback,main}%
\tikzset{
  fig/.style={x=1mm, y=1mm, font=\figsf\fontsize{6.6}{7.8}\selectfont, text=fink},
  figsmall/.style={font=\figsf\fontsize{5.9}{6.9}\selectfont},
  fighead/.style={font=\figsf\bfseries\fontsize{7.2}{8.4}\selectfont, text=fslate},
  figgloss/.style={font=\figsf\itshape\fontsize{6.1}{7.1}\selectfont, text=fmute},
  figpanel/.style={draw=fedge, fill=white, rounded corners=2.2mm, line width=0.4pt},
  figarrow/.style={-{Stealth[length=1.6mm, width=1.2mm]}, line width=0.6pt},
}%
\fi
% make_thesis_assets.py copies it to generated/. Every line ends in % because figure files are read
% inside boxes, where a stray space would widen the figure.
\def\thesisfigstyle{}%
\DeclareFontFamily{TU}{pipeheros}{}%
\DeclareFontShape{TU}{pipeheros}{m}{n}{<->"[texgyreheros-regular.otf]:mapping=tex-text"}{}%
\DeclareFontShape{TU}{pipeheros}{b}{n}{<->"[texgyreheros-bold.otf]:mapping=tex-text"}{}%
\DeclareFontShape{TU}{pipeheros}{m}{it}{<->"[texgyreheros-italic.otf]:mapping=tex-text"}{}%
\DeclareFontShape{TU}{pipeheros}{b}{it}{<->"[texgyreheros-bolditalic.otf]:mapping=tex-text"}{}%
\DeclareFontShape{TU}{pipeheros}{bx}{n}{<->ssub*pipeheros/b/n}{}%
\DeclareFontShape{TU}{pipeheros}{bx}{it}{<->ssub*pipeheros/b/it}{}%
\newcommand{\figsf}{\fontfamily{pipeheros}}%
% palette: honest green and deceptive red are also the steering figure's colours
\definecolor{fhonest}{HTML}{1B8A5A}%
\definecolor{fhonestbg}{HTML}{DDF0E6}%
\definecolor{fdecept}{HTML}{C8423B}%
\definecolor{fdeceptbg}{HTML}{F7E0DE}%
\definecolor{famber}{HTML}{B7791F}%
\definecolor{fambersbg}{HTML}{F8EBD3}%
\definecolor{fslate}{HTML}{2E3A4F}%
\definecolor{fslatebg}{HTML}{E4E7EC}%
\definecolor{fcanvas}{HTML}{F4F2EE}%
\definecolor{fedge}{HTML}{D8D3CA}%
\definecolor{fink}{HTML}{1F2328}%
\definecolor{fmute}{HTML}{6B6B6B}%
\pgfdeclarelayer{figbase}%
\pgfdeclarelayer{figback}%
\pgfsetlayers{figbase,figback,main}%
\tikzset{
  fig/.style={x=1mm, y=1mm, font=\figsf\fontsize{6.6}{7.8}\selectfont, text=fink},
  figsmall/.style={font=\figsf\fontsize{5.9}{6.9}\selectfont},
  fighead/.style={font=\figsf\bfseries\fontsize{7.2}{8.4}\selectfont, text=fslate},
  figgloss/.style={font=\figsf\itshape\fontsize{6.1}{7.1}\selectfont, text=fmute},
  figpanel/.style={draw=fedge, fill=white, rounded corners=2.2mm, line width=0.4pt},
  figarrow/.style={-{Stealth[length=1.6mm, width=1.2mm]}, line width=0.6pt},
}%
\fi%
\begin{tikzpicture}[fig,
    word/.style={circle, fill=fink, inner sep=0pt, minimum size=1.5mm},
    wlab/.style={figsmall, inner sep=1.5mm},
    plane/.style={fill=fcanvas, draw=fedge, line width=0.5pt},
    eq/.style={draw=fslate!45, fill=white, rounded corners=1.2mm, inner sep=1.3mm, align=center}]
% \plane{x0}: a tilted plane with a light grid, 46 mm wide, from y = 26 to 62
\newcommand{\plane}[1]{%
  \filldraw[plane] (#1,26) -- (#1+46,26) -- (#1+52,62) -- (#1+6,62) -- cycle;
  \foreach \t in {0.2,0.4,0.6,0.8} {
    \draw[fedge, line width=0.3pt] ($(#1,26)!\t!(#1+6,62)$) -- ($(#1+46,26)!\t!(#1+52,62)$);
    \draw[fedge, line width=0.3pt] ($(#1,26)!\t!(#1+46,26)$) -- ($(#1+6,62)!\t!(#1+52,62)$);}}

% ---------------------------------------------------------------- left: measure
\draw[figpanel] (0,0) rectangle (62,80);
\node[fighead, anchor=north west] at (3,78) {Unembedding space $\Gamma$ (words)};
\node[figgloss, anchor=north west] at (3,73.4) {measure: read a concept by projection};
\plane{5}
\node[word] (man) at (16,31) {};   \node[wlab, below left] at (man) {man};
\node[word] (woman) at (27,38) {}; \node[wlab, right] at (woman) {woman};
\node[word] (king) at (18,45) {};  \node[wlab, left] at (king) {king};
\node[word] (queen) at (29,52) {}; \node[wlab, above] at (queen) {queen};
\node[word] (roi) at (40,45) {};   \node[wlab, below] at (roi) {roi};
\node[word] (reine) at (51,52) {}; \node[wlab, above] at (reine) {reine};
\foreach \a/\b in {man/woman, king/queen, roi/reine}
  \draw[figarrow, fdecept, shorten >=0.9mm, shorten <=0.9mm] (\a) -- (\b);
\draw[figarrow, fslate, dashed, shorten >=0.9mm, shorten <=0.9mm] (king) -- (roi);
\node[figsmall, text=fslate] at (29,47.4) {$\bar\gamma_Z$};
\node[figsmall, text=fdecept, anchor=west] at (31,31) {$\bar\gamma_W$: male$\Rightarrow$female};
\node[eq, anchor=south] at (31,4) {$\operatorname{logit} P(W{=}1\mid\lambda)=\alpha\,\lambda^{\top}\bar\gamma_W$\\[0.3ex]
  {\figsf\itshape\fontsize{6}{7}\selectfont\color{fmute} probing reads $W$ along $\bar\gamma_W$}};
\node[figsmall, text=fmute, align=center] at (31,20.5) {three word pairs share one direction};

% ---------------------------------------------------------------- centre: the causal inner product
\draw[fslatebg, fill=fslatebg] (75,55) .. controls (81,45) .. (75,35) .. controls (69,45) .. (75,55);
\draw[fslate, line width=0.5pt] (75,55) .. controls (81,45) .. (75,35) .. controls (69,45) .. (75,55);
\node[figsmall, text=fslate, align=center] at (75,45) {causal\\[-0.4ex]inner\\[-0.4ex]product};
\draw[figarrow, fslate] (62.5,45) -- (71.6,45);
\draw[figarrow, fslate] (78.4,45) -- (87.5,45);
\node[figsmall, align=center, text=fslate] at (75,27.5)
  {$\bar\lambda_W=$\\$\mathrm{Cov}(\gamma)^{-1}\bar\gamma_W$};
\node[figgloss, align=center] at (75,64) {$Z\perp W$\\[-0.3ex]under it};

% ---------------------------------------------------------------- right: steer
\draw[figpanel] (88,0) rectangle (150,80);
\node[fighead, anchor=north west] at (91,78) {Embedding space $\Lambda$ (contexts)};
\node[figgloss, anchor=north west] at (91,73.4) {steer: write a concept by addition};
\plane{93}
\fill[fmute!25] (104,33) ellipse (4.2 and 2.6);
\node[word] (lx) at (104,33) {};
\node[wlab, below] at (104,30.2) {$\lambda(x)$: ``the \dots{} was crowned''};
\node[word] (lxw) at (128,53) {};
\node[wlab, above] at (lxw) {$\lambda(x)+c\,\bar\lambda_W$};
\draw[figarrow, fdecept, line width=0.9pt, shorten >=0.9mm, shorten <=0.9mm] (lx) -- (lxw)
  node[midway, figsmall, text=fdecept, anchor=north west] {$+\,c\,\bar\lambda_W$};
% next-token bars before and after (schematic)
\foreach \x/\k/\q/\lab in {95/9/2/before, 125/2/9/after} {
  \fill[fslate!55] (\x,10) rectangle ++(4,\k);
  \fill[fdecept!80] (\x+8,10) rectangle ++(4,\q);
  \draw[fmute, line width=0.3pt] (\x-1.5,10) -- (\x+13.5,10);
  \node[figsmall, anchor=north] at (\x+2,9.8) {king};
  \node[figsmall, anchor=north] at (\x+10,9.8) {queen};
  \node[figsmall, text=fmute, anchor=south] at (\x+6,20.2) {\lab};}
\draw[figarrow, fdecept] (110.5,15) -- (122.5,15);
\node[figgloss, anchor=north] at (119,4.8) {next-token probability (schematic)};
\pgfresetboundingbox
\path (0,0) rectangle (150,80);
\end{tikzpicture}%

\caption{The linear representation hypothesis as formalised by \textcite{park2024linear}.}
\label{fig:lrh}
\end{figure}

\emph{Why linearity arises, and where it fails.} Next-token prediction trained by gradient descent
promotes linear representations of concepts \cite{jiang2024origins}. For truth in particular, true
statements tend to co-occur with true ones in text, so a latent ``this text is truthful'' variable
lowers the loss, and a toy transformer learns to separate true from false linearly
\cite{ravfogel2025emergence}. Linear directions both read and write: a single difference-of-means
direction mediates refusal in thirteen chat models of up to 72B parameters, including Llama-3-70B
\cite{arditi2024refusal}, and differences averaged over contrast pairs steer behaviours such as
sycophancy \cite{rimsky2024steering}. But the hypothesis holds only in a weak form. Linearity depends
on the variable: the board state of a transformer trained on Othello looked nonlinear in terms of black
and white \cite{li2023emergent}, yet is 99.5\% linearly decodable in terms of ``mine'' and ``yours''
\cite{nanda2023emergent}. And some features are irreducibly multi-dimensional, such as the circles on
which days of the week and months lie \cite{engels2025not}.

\emph{Implications for detection.} Detection relies on the measurement reading and steering on the
intervention reading, and \cref{eq:causal-ip} says their directions need not coincide,
as the truth probes of \cref{sec:review-truth} illustrate. No linear classifier can beat a constant predictor
exactly when the class-conditional means coincide \cite{belrose2023leace}, so the difference of means
carries all of a binary concept's linearly available information. A persona contrast
(\cref{eq:contrast}) is a difference of means and so a candidate for both reading and writing, while a
supervised probe is fitted only to read; \cref{sec:steering} tests the difference.

\subsection{Probing techniques}
\label{sec:review-probes}

\emph{What a probe measures.} Probe accuracy is hard to interpret on its own. An expressive probe can
memorise: on part-of-speech tags read from ELMo, a multi-layer perceptron matched a linear probe's
accuracy but also scored 92.8\% on a \emph{control task} of random labels, against 71.2\% for the
linear probe \cite{hewitt2019designing}. Decodable information need not be used, since removing it often
leaves behaviour unchanged \cite{elazar2021amnesic}, and probing datasets are imperfect proxies for
the task \cite{belinkov2022probing}. For lie detection the information is present in some form; the
question is whether a simple readout extracts the intended variable or a correlate of it.

\emph{Families of probe.} \Cref{tab:probes} lists the main families. Supervised linear readouts
differ in how they choose the direction (\cref{sec:prelim-probes}): logistic regression tilts away from
a concept's direction to cancel correlated features, while the difference of means points at it
\cite{marks2023geometry}. Label-free readouts replace labels with structure and cannot fit a dataset's
artefacts, but they pick up whatever most separates their contrast: unsupervised methods ``seem to
discover whatever feature of the activations is most prominent'' \cite{farquhar2023challenges}.
Sparse-autoencoder features brought no consistent gain \cite{kantamneni2025sparse}, and how per-token
scores are pooled matters: the last token did worst for high-stakes detection
\cite{mckenzie2025highstakes}.

\begin{table}[t]
\centering
\small
\caption[Families of probing technique]{Families of probing technique and the weakness documented for
each.}
\label{tab:probes}
\begin{tabularx}{\linewidth}{@{}>{\raggedright\arraybackslash}p{3.0cm}>{\raggedright\arraybackslash}p{4.6cm}>{\raggedright\arraybackslash}X@{}}
\toprule
Family & Examples & Documented weakness \\
\midrule
Logistic regression & Mean Probe \cite{kretschmar2025liars}, deception probes
\cite{goldowskydill2025detecting} & Fits dataset artefacts \cite{levinstein2025still} \\
Difference of means, LDA & mass-mean probe \cite{marks2023geometry}, steering vectors
\cite{rimsky2024steering} & Absorbs any systematic difference between classes \\
Unsupervised & CCS \cite{burns2023discovering}, LAT \cite{zou2023representation} & Finds the most
prominent feature \cite{farquhar2023challenges} \\
Contrast prompts & persona vectors \cite{chen2025persona}, persona coordinates
\cite{mahadik2026persona} & Reads whatever the prompts differ in \\
Sparse autoencoder & SAE probes \cite{kantamneni2025sparse} & No consistent gain; worse under shift \\
Follow-up question & ``Was your last response deceptive?'' \cite{goldowskydill2025detecting} & Extra
tokens per response \\
Token pooling & mean, max, softmax, attention \cite{mckenzie2025highstakes} & Mean dilutes partial
lies; long contexts \cite{kramar2026production} \\
\bottomrule
\end{tabularx}
\end{table}

\subsection{Linear representations of truth}
\label{sec:review-truth}

\emph{First probes and their critics.} A classifier on a model's hidden states separates true from
false statements, with accuracy that varies widely across held-out topics \cite{azaria2023internal};
truth directions can also be found without labels, from the consistency of a model's answers to a
statement and its negation \cite{burns2023discovering}. But any binary feature satisfies that
consistency objective, and in experiments unsupervised methods tracked whether a prompt ended in
``banana'', or what a simulated character thought, instead of the label of interest
\cite{farquhar2023challenges}.

\emph{The geometry of truth.} \textcite{marks2023geometry} give the strongest evidence for a linear
truth direction. On curated statements such as ``The city of Tokyo is in Japan'', the top principal
components of LLaMA-2-70B's residual stream separate true from false; probes trained on one dataset
generalise to others at 13B and 70B; and adding the probe direction shifts the model's true/false
judgements on a dataset no probe saw. Two controls bear on spurious features. Text probability
correlates with truth on some datasets and against it on negated ones, so a probe trained to separate
likely from unlikely text is worse than chance there. And the difference of means generalises about as
well as logistic regression while being far more causal: for LLaMA-2-13B trained on cities, it moves
the model's judgement of false statements 77\% of the way to true ones, logistic regression 13\%.

\emph{Still no lie detector.} \textcite{levinstein2025still} showed the failure such controls
anticipate: probes trained on affirmative statements scored 0.37--0.41 on the same facts negated, and
CCS probes learned which member of each pair was negated rather than which was true. The probes, they
conclude, learn ``some other feature that correlates well with truth on the training sets but that does
not correlate with truth in even mildly more general contexts''.

\emph{Subspaces and hidden knowledge.} Truth occupies at least two dimensions, a general direction and
a polarity direction that flips with negation; probes trained on affirmative statements learn a
mixture of the two \cite{buerger2024truth}. Diversity of training data helps more than volume
\cite{liu2024universal}, yet error-detection probes ``fail to generalize across datasets'', and much of
the transfer that remains is matched by the output logits \cite{orgad2025llms}. Models also represent
truths they do not say: probe accuracy can exceed generation accuracy by 40 points
\cite{li2023inference}, and when LLaMA-2-70B-chat is instructed to lie, truth probes trained on
honest and on lying prompts become anti-parallel after about layer~23 \cite{campbell2023localizing}. A
denial of knowledge differs from all of these: it asserts no false fact for a truth probe to find.

\subsection{The truthfulness spectrum hypothesis}
\label{sec:review-tsh}

\textcite{ying2026truthfulness} argue that a probe's failure to transfer does not show that no
domain-general truth direction exists, and propose the \emph{truthfulness spectrum hypothesis}:
``LLMs encode truthfulness along a spectrum of generality, with directions at varying levels of
generality coexisting in the representational space''. Their setting is close to this thesis:
Llama-3.3-70B-Instruct, and difference-of-means, logistic-regression and LDA probes on activations
averaged over output tokens, scored by AUROC. The data cover five
truth types (definitional, empirical, logical, fictional and ethical) and two stress tests:
\emph{sycophantic lying}, collected on-policy as changes of answer towards a user's stated answer, and
\emph{expectation-inverted lying}, in which the model is told to ``Pretend you're a dishonest
person'', so that the true statements are the unexpected ones.

Four results support the hypothesis. (1)~Probes transfer well across truth types but not to or from the
two kinds of lie; a probe trained on expectation-inverted lying scores ``well below chance
(AUROC$\approx$0.28)'' on the truth types. (2)~A probe trained on all domains jointly nevertheless
performs well everywhere: pairwise failure ``reflects incomplete recovery of [the] general direction,
not its absence''. (3)~Transfer from domain $A$ to domain $B$ is predicted by the Mahalanobis cosine
similarity of the two probes' weights,
\begin{equation}
  \cos_{\Sigma}(w_A, w_B) \;=\; \frac{w_A^{\top}\Sigma\, w_B}
  {\sqrt{w_A^{\top}\Sigma\, w_A}\,\sqrt{w_B^{\top}\Sigma\, w_B}} ,
  \label{eq:mahalanobis}
\end{equation}
with $\Sigma$ the covariance of the test data, at $R^2 = 0.98$ against $0.56$ for the ordinary cosine,
which fails because ``the effective dimensionality of our data is fewer than 100 in an
8192-dimensional space''. (4)~Iterative nullspace projection and LEACE erasure isolate directions of
every generality: domain-general, domain-specific, and shared by subsets of domains. Post-training
moves sycophancy away from the other types (probes trained on the truth types reach 0.74 AUROC on
sycophantic lying in base Llama models but 0.62 (70B) and 0.59 (8B) in chat models), and in a causal
test the domain-specific directions increased the model's preference for correct answers while the
general direction decreased it.

The hypothesis offers three lessons for this thesis. A probe trained on one Liars' Bench dataset can
transfer poorly even if a general deception direction exists, because it mixes that direction with
dataset-specific ones. Expectation-inverted lying, which reverses a surface correlation and inverts the
probe trained on it, has the shape of \hpkr: its lies are denials, while every published deceptive
persona lies by controlling disclosure, never by denial. And training on many domains has a label-free analogue in widening a persona
bank, though ``coverage is never guaranteed''.

\subsection{Probes for deception}
\label{sec:review-deception}

\textcite{goldowskydill2025detecting} trained logistic-regression probes on layer~22 of
Llama-3.3-70B-Instruct, on instruction pairs or role-play, and averaged their per-token scores. On the
model's own deception in insider trading and sandbagging they reached AUROCs of 0.96--0.999 and caught
95--99\% of deceptive responses at 1\% false positives on Alpaca, yet the authors judge this
``insufficient as a robust defence''. The probe fires more on honest responses to the evaluation
prompts than on control prompts, because it responds to whether ``the scenario [is] related to
deception'' and to morality (honest confessions of misdeeds score as dishonest); averaging misses
partly deceptive responses; and two layers later, recall on insider trading drops to nearly zero.

\textcite{mckenzie2025highstakes} show a success for a property visible in the text. Probes on layer~31
of Llama-3.3-70B, trained on synthetic data, detected high-stakes interactions in six out-of-distribution datasets with a mean AUROC above 0.91,
``comparable to that of prompted or finetuned medium-sized LLM monitors'' at a millionth of the cost.
Yet at 1\% false positives the best probe caught 43\% of cases against 55\% for the best language-model
baseline, and probes still reacted to surface cues, flagging a harmless question ``due to the word
`emergency' ''. For deception, confidently labelled examples of strategic lies are scarce
\cite{smith2025difficulties}. Probes can also be evaded: adversarial suffixes can cut the recall of
harmfulness probes from 100\% to 0\% \cite{bailey2026obfuscated}, and models trained to evade probes for
benign concepts evade deception probes they have never seen \cite{mcguinness2025neural}.

\subsection{Black-box lie detection}
\label{sec:review-blackbox}

Black-box detectors use only inputs and outputs. \textcite{pacchiardi2023catch} ask a fixed set of
unrelated follow-up questions after a suspected lie and classify the answers, and Liars' Bench adds an
LLM judge that reads the transcript and a self-evaluation in which the model is asked whether it lied
\cite{kretschmar2025liars}. The judge is the strongest overall but fails on \hpkr, one of the settings
where ``it's not possible to determine whether the model lied from the transcript alone''. EleutherAI's entry to the
Aletheia's Quest competition combined a black-box auditor with a white-box self-evaluation channel read
in a third-person auditing context \cite{birardi2026aletheia}. Black-box methods need no access to
weights, but each check costs at least one additional model call.

\subsection{Liars' Bench}
\label{sec:review-lb}

Liars' Bench \cite{kretschmar2025liars} collects \V{lb-examples} on-policy lies and honest responses
across seven datasets and four models, classified by the object of belief and the reason for lying
(\cref{sec:prelim-lies}). \Cref{tab:datasets} lists the five datasets that Llama-3.3-70B-Instruct
generated without fine-tuning, which this thesis uses. The benchmark evaluates a \emph{Mean Probe}
trained on RepE instruction pairs, a \emph{Follow-Up Probe} read at a prefilled ``No'' after the
question ``Was your last response deceptive?'', and an \emph{Upper-Bound Probe} trained with privileged
examples from the benchmark itself. Most methods perform near chance on the harm-pressure datasets. On
\hpkr, built from the hazardous-knowledge questions of WMDP \cite{li2024wmdp}, the judge's recall is
almost zero and the Mean and Follow-Up Probes invert (\V{lb-mean-auroc-hpkr} and
\V{lb-followup-auroc-hpkr} AUROC). The unrelated-question classifier ranks \hpkr's lies well (an AUROC
of 0.89 is reported for Llama-3.3-70B), but it needs 39 extra queries per response and is described as
typically near chance in balanced accuracy on these datasets. Only the Upper-Bound Probe, trained with
labelled \hpkr{} examples, detects them reliably, at \V{lb-upper-balacc-hpkr} balanced accuracy. The
authors attribute the Mean Probe's failure to features ``only superficially correlated with lying on
the training dataset''. A later study found ``sub-random AUROC performance (e.g., 0.12--0.16) on the
HP-KR dataset across nearly all probes'' in two other models, and suggested that probes trained on
direct lies ``latch onto a feature direction that maps inversely'' to \hpkr's denials
\cite{moustafa2026beyond}.

\begin{table}[t]
\centering
\small
\caption{Datasets. Object of belief and reason for lying follow Liars' Bench's taxonomy. Counts are
the Llama-3.3-70B-Instruct examples we score. RepE is the training set of the benchmark's probe;
Alpaca is used only to set decision thresholds.}
\label{tab:datasets}
\fittable{% Generated by make_numbers.py
\begin{tabular}{llllrr}
\toprule
Dataset & Short & Object of belief & Reason for lying & Examples & Lies \\
\midrule
Convincing-Game & CG & self-knowledge & encouraged & 267 & 167 \\
Harm-Pressure Choice & HP-C & world-knowledge & inherent & 2,134 & 134 \\
Harm-Pressure Knowledge Report & HP-KR & self-knowledge & inherent & 2,139 & 139 \\
Insider-Trading & IT & self-knowledge & encouraged & 1,080 & 423 \\
Instructed-Deception & ID & world-knowledge & encouraged & 5,432 & 3,045 \\
\midrule
RepE instruction pairs & RepE & world-knowledge & instructed & 550 & 275 \\
Alpaca (control) & n/a & n/a & none & 2,000 & 0 \\
\bottomrule
\end{tabular}}
\end{table}

\Cref{fig:scenarios} in \cref{sec:data} shows one lie by Llama-3.3-70B-Instruct from each of these
five datasets, placed on the benchmark's two axes.


\subsection{Personas as directions}
\label{sec:review-personas}

A line of work treats honesty as a property of a speaker the model simulates. A model trained on
text from many authors must infer a context's likely author and represent that author's beliefs and
intentions \cite{andreas2022language}; a dialogue agent says something false ``deliberately'' when ``it
is role-playing a deceptive character'' \cite{shanahan2023role}; and a model enacts the Assistant with
the same concepts it uses for human and fictional characters \cite{marks2026persona}. A probe on a
question alone predicts whether the answer will be truthful, as if the model inferred a truthful or
untruthful persona \cite{joshi2024personas}, and a ``toxic persona'' feature controls emergent
misalignment \cite{wang2026persona}.

\textcite{chen2025persona} introduced \emph{persona vectors}: activation directions for character traits
such as sycophancy, extracted automatically from a natural-language description of the trait and used
to monitor and steer it. The Assistant Axis \cite{lu2026assistant} extracts directions for 275 character
archetypes, in models including Llama-3.3-70B, and finds that the leading component of this persona
space measures how far the model sits from its default assistant persona.
\textcite{mahadik2026persona} turned this machinery into a monitor for deception and sycophancy. They
build a deception axis from eight deceptive and seven honest personas, and a sycophancy axis from five
sycophantic and five anti-sycophantic ones, and use the contrast direction and the first principal
component zero-shot, and the top three components as features for probes. On Llama~3.2-3B and
Llama~3-8B, these persona coordinates transfer across ten deception and sycophancy datasets
(including Liars' Bench's Convincing-Game, Instructed-Deception and Harm-Pressure Choice, but not \hpkr)
better than probes on raw activations. They also build a \emph{unified} harmful-behaviour axis by
pooling the deceptive and sycophantic personas. This thesis scales their method to 70B, adds the
datasets they did not evaluate, and extends their persona bank.

The same literature marks the limits of the approach. A role-played lie may not be the model's own
deception \cite{smith2025difficulties}; roles differ in style as well as in honesty, and
\textcite{joshi2024personas} suspect their own persona signal of encoding style; and a model that
switches to an ignorant persona deceives although ``no persona is ever lying'' \cite{marks2026persona},
which is the shape of a denial of knowledge.

\subsection{Why probes fail to transfer}
\label{sec:review-transfer}

Across this literature, a probe that is nearly perfect in distribution can be near or below chance
elsewhere. This thesis finds the same at 70B: probes trained on one Liars' Bench dataset lose much of their
in-distribution AUROC on the others, and one train--test pair falls far below chance
(\cref{sec:supervised}). The literature points to six mechanisms, the central one illustrated in
\cref{fig:spurious}.

\emph{1. Underspecification.} Every dataset in \cref{tab:datasets} has fewer examples than the 8192
dimensions of the residual stream, and $N \le D$ points in general position can be split by a hyperplane
in all $2^N$ ways \cite{cover1965geometrical}. Real activations concentrate their variance in fewer than
a hundred directions \cite{ying2026truthfulness}, but many directions still fit a training set equally
well. Such predictors ``can behave very differently in deployment domains''
\cite{damour2022underspecification}, and logistic regression chooses among them by the margin of the
particular training set \cite{soudry2018implicit}.

\emph{2. Shortcuts.} Discriminative learning uses ``any feature that is sufficient to reliably
discriminate on a given dataset'' \cite{geirhos2020shortcut}. Probes have latched
onto negation and polarity \cite{levinstein2025still,buerger2024truth}, text probability
\cite{marks2023geometry}, topic and required skill \cite{orgad2025llms}, deception-related scenarios and
morality \cite{goldowskydill2025detecting}, single words such as ``emergency''
\cite{mckenzie2025highstakes},
style \cite{kumar2026pressure} and a simulated character's belief \cite{farquhar2023challenges}.

\emph{3. Why the shortcut wins.} Gradient-trained linear classifiers use a spurious feature even when
an invariant one predicts the label perfectly, because it lets them separate the data with a smaller
norm \cite{nagarajan2021understanding}; the most linearly decodable feature is preferred, and ``an
easier, weakly predictive feature can suppress a more strongly predictive, but more difficult one''
\cite{hermann2020what}. For deception, the directions a probe weights heavily ``appear to encode
source-specific surface features'', while the directions that transfer encode the contrast more
abstractly \cite{yoo2026subspace}.

\Cref{fig:spurious} draws mechanisms 2 and 3 as a schematic. In panel (a), the lies of dataset $A$
differ from its honest responses along a feature $u$ that all lies share and, more strongly, along a
feature $s$ specific to $A$, and the maximum-margin probe $w_A$ leans on $s$. In panel (b), dataset $B$
reverses the relation of $s$ to the label, so $w_A$ ranks lies below honest responses, while a fixed
direction $d$ aligned with $u$ still ranks them correctly.

\begin{figure}[t]
\centering
% How a probe inherits a dataset-specific feature (fig:spurious), schematic. Two datasets whose lies
% share a feature u; in A the lies also differ strongly along s, and in B that relation is reversed.
% Under each scatter, the score densities of honest responses and lies projected on the probe w_A and
% on the fixed direction d. Geometry: honest at (-0.75, -1.3 sgn), lies at (0.75, 1.3 sgn), w_A along
% (0.5, 0.9), d along (1, 0), so on w_A the class means are +-1.5 in A and -+0.77 in B.
\ifdefined\thesisfigstyle\else% Shared style of the thesis's TikZ figures: the TeX Gyre Heros face (declared for XeTeX's TU
% encoding), the thesis palette and two drawing layers. Each figure file loads it once through
%   \ifdefined\thesisfigstyle\else% Shared style of the thesis's TikZ figures: the TeX Gyre Heros face (declared for XeTeX's TU
% encoding), the thesis palette and two drawing layers. Each figure file loads it once through
%   \ifdefined\thesisfigstyle\else\input{generated/fig_style}\fi
% make_thesis_assets.py copies it to generated/. Every line ends in % because figure files are read
% inside boxes, where a stray space would widen the figure.
\def\thesisfigstyle{}%
\DeclareFontFamily{TU}{pipeheros}{}%
\DeclareFontShape{TU}{pipeheros}{m}{n}{<->"[texgyreheros-regular.otf]:mapping=tex-text"}{}%
\DeclareFontShape{TU}{pipeheros}{b}{n}{<->"[texgyreheros-bold.otf]:mapping=tex-text"}{}%
\DeclareFontShape{TU}{pipeheros}{m}{it}{<->"[texgyreheros-italic.otf]:mapping=tex-text"}{}%
\DeclareFontShape{TU}{pipeheros}{b}{it}{<->"[texgyreheros-bolditalic.otf]:mapping=tex-text"}{}%
\DeclareFontShape{TU}{pipeheros}{bx}{n}{<->ssub*pipeheros/b/n}{}%
\DeclareFontShape{TU}{pipeheros}{bx}{it}{<->ssub*pipeheros/b/it}{}%
\newcommand{\figsf}{\fontfamily{pipeheros}}%
% palette: honest green and deceptive red are also the steering figure's colours
\definecolor{fhonest}{HTML}{1B8A5A}%
\definecolor{fhonestbg}{HTML}{DDF0E6}%
\definecolor{fdecept}{HTML}{C8423B}%
\definecolor{fdeceptbg}{HTML}{F7E0DE}%
\definecolor{famber}{HTML}{B7791F}%
\definecolor{fambersbg}{HTML}{F8EBD3}%
\definecolor{fslate}{HTML}{2E3A4F}%
\definecolor{fslatebg}{HTML}{E4E7EC}%
\definecolor{fcanvas}{HTML}{F4F2EE}%
\definecolor{fedge}{HTML}{D8D3CA}%
\definecolor{fink}{HTML}{1F2328}%
\definecolor{fmute}{HTML}{6B6B6B}%
\pgfdeclarelayer{figbase}%
\pgfdeclarelayer{figback}%
\pgfsetlayers{figbase,figback,main}%
\tikzset{
  fig/.style={x=1mm, y=1mm, font=\figsf\fontsize{6.6}{7.8}\selectfont, text=fink},
  figsmall/.style={font=\figsf\fontsize{5.9}{6.9}\selectfont},
  fighead/.style={font=\figsf\bfseries\fontsize{7.2}{8.4}\selectfont, text=fslate},
  figgloss/.style={font=\figsf\itshape\fontsize{6.1}{7.1}\selectfont, text=fmute},
  figpanel/.style={draw=fedge, fill=white, rounded corners=2.2mm, line width=0.4pt},
  figarrow/.style={-{Stealth[length=1.6mm, width=1.2mm]}, line width=0.6pt},
}%
\fi
% make_thesis_assets.py copies it to generated/. Every line ends in % because figure files are read
% inside boxes, where a stray space would widen the figure.
\def\thesisfigstyle{}%
\DeclareFontFamily{TU}{pipeheros}{}%
\DeclareFontShape{TU}{pipeheros}{m}{n}{<->"[texgyreheros-regular.otf]:mapping=tex-text"}{}%
\DeclareFontShape{TU}{pipeheros}{b}{n}{<->"[texgyreheros-bold.otf]:mapping=tex-text"}{}%
\DeclareFontShape{TU}{pipeheros}{m}{it}{<->"[texgyreheros-italic.otf]:mapping=tex-text"}{}%
\DeclareFontShape{TU}{pipeheros}{b}{it}{<->"[texgyreheros-bolditalic.otf]:mapping=tex-text"}{}%
\DeclareFontShape{TU}{pipeheros}{bx}{n}{<->ssub*pipeheros/b/n}{}%
\DeclareFontShape{TU}{pipeheros}{bx}{it}{<->ssub*pipeheros/b/it}{}%
\newcommand{\figsf}{\fontfamily{pipeheros}}%
% palette: honest green and deceptive red are also the steering figure's colours
\definecolor{fhonest}{HTML}{1B8A5A}%
\definecolor{fhonestbg}{HTML}{DDF0E6}%
\definecolor{fdecept}{HTML}{C8423B}%
\definecolor{fdeceptbg}{HTML}{F7E0DE}%
\definecolor{famber}{HTML}{B7791F}%
\definecolor{fambersbg}{HTML}{F8EBD3}%
\definecolor{fslate}{HTML}{2E3A4F}%
\definecolor{fslatebg}{HTML}{E4E7EC}%
\definecolor{fcanvas}{HTML}{F4F2EE}%
\definecolor{fedge}{HTML}{D8D3CA}%
\definecolor{fink}{HTML}{1F2328}%
\definecolor{fmute}{HTML}{6B6B6B}%
\pgfdeclarelayer{figbase}%
\pgfdeclarelayer{figback}%
\pgfsetlayers{figbase,figback,main}%
\tikzset{
  fig/.style={x=1mm, y=1mm, font=\figsf\fontsize{6.6}{7.8}\selectfont, text=fink},
  figsmall/.style={font=\figsf\fontsize{5.9}{6.9}\selectfont},
  fighead/.style={font=\figsf\bfseries\fontsize{7.2}{8.4}\selectfont, text=fslate},
  figgloss/.style={font=\figsf\itshape\fontsize{6.1}{7.1}\selectfont, text=fmute},
  figpanel/.style={draw=fedge, fill=white, rounded corners=2.2mm, line width=0.4pt},
  figarrow/.style={-{Stealth[length=1.6mm, width=1.2mm]}, line width=0.6pt},
}%
\fi%
\begin{tikzpicture}[fig,
    hon/.style={circle, draw=fhonest, fill=white, line width=0.5pt, inner sep=0pt, minimum size=1.7mm},
    lie/.style={rectangle, draw=fdecept, fill=fdecept!80, line width=0.4pt, inner sep=0pt,
                minimum size=1.5mm},
    ax/.style={-{Stealth[length=1.2mm]}, draw=fmute, line width=0.4pt}]
% \strip{x centre}{y baseline}{honest mean}{lie mean}{label}{verdict}: two score densities
\newcommand{\strip}[6]{%
  \draw[fmute, line width=0.4pt] (#1-25,#2) -- (#1+25,#2);
  \draw[fmute, line width=0.3pt] (#1,#2-0.8) -- (#1,#2+0.4);
  \fill[fhonestbg, draw=fhonest, line width=0.5pt, fill opacity=0.85]
    plot[domain=-25:25, samples=70, smooth] ({#1+\x},{#2+8*exp(-((\x/10-(#3))^2)/(2*0.33^2))}) -- cycle;
  \fill[fdeceptbg, draw=fdecept, line width=0.5pt, fill opacity=0.75]
    plot[domain=-25:25, samples=70, smooth] ({#1+\x},{#2+8*exp(-((\x/10-(#4))^2)/(2*0.33^2))}) -- cycle;
  \node[figsmall, anchor=north west, text=fmute, inner sep=0.5pt] at (#1-25,#2-0.9) {#5};
  \node[figsmall, anchor=north east, inner sep=0.5pt] at (#1+25,#2-0.9) {#6};}

\foreach \sgn/\cx/\head/\sub in {1/37/{(a) Training dataset $A$}/{the probe leans on $s$},
                                -1/112/{(b) Test dataset $B$}/{the relation of $s$ is reversed}} {
  \draw[figpanel] (\cx-35,0) rectangle (\cx+35,103);
  \node[fighead, anchor=north west] at (\cx-32,100.5) {\head};
  \node[figgloss, anchor=north west] at (\cx-32,96) {\sub};
  \begin{scope}[shift={(\cx,62)}, x=10mm, y=10mm]
    % soft class clouds
    \fill[fhonestbg, opacity=0.9] (-0.75,{-1.3*\sgn}) ellipse (0.62 and 0.55);
    \fill[fdeceptbg, opacity=0.9] (0.75,{1.3*\sgn}) ellipse (0.62 and 0.55);
    \draw[ax] (-2.5,0) -- (2.5,0);
    \draw[ax] (0,-2.35) -- (0,2.45);
    \node[figsmall, text=fmute, anchor=north east, align=right] at (2.5,-0.12) {shared lie\\[-0.4ex]feature $u$};
    \node[figsmall, text=fmute, anchor=south west] at (0.08,2.3) {dataset feature $s$};
    \foreach \dx/\dy in {0/0, .28/.18, -.22/.3, .15/-.27, -.32/-.12, .36/-.06, -.1/.42, .05/-.4, .22/.34, -.26/-.33}{
      \node[hon] at ({-0.75+\dx},{\sgn*(-1.3)+\dy}) {};
      \node[lie] at ({0.75-\dx},{\sgn*1.3-\dy}) {};
    }
    % the probe's boundary and the two directions
    \draw[fslate, dashed, line width=0.6pt] (-2.1,1.18) -- (2.1,-1.18);
    \draw[figarrow, fslate] (0,0) -- (0.62,1.12);
    \node[figsmall, text=fslate, anchor=east] at (0.26,0.62) {$w_A$};
    \draw[figarrow, famber, line width=0.9pt] (0,0) -- (1.75,0);
    \node[figsmall, text=famber, anchor=south] at (1.6,0.05) {$d$};
  \end{scope}
}
% score strips: A separates on both; B inverts on w_A and still separates on d
\strip{37}{22}{-1.5}{1.5}{score on $w_A$}{\textcolor{fhonest}{$\checkmark$ lies score higher}}
\strip{37}{5}{-0.75}{0.75}{score on $d$}{\textcolor{fhonest}{$\checkmark$ lies score higher}}
\strip{112}{22}{0.77}{-0.77}{score on $w_A$}{\textcolor{fdecept}{$\times$ inverted: lies score lower}}
\strip{112}{5}{-0.75}{0.75}{score on $d$}{\textcolor{fhonest}{$\checkmark$ still in order}}
% key
\node[hon] at (6,33.5) {};
\node[figsmall, anchor=west] at (7.4,33.5) {honest};
\node[lie] at (20,33.5) {};
\node[figsmall, anchor=west] at (21.4,33.5) {lie};
\node[figsmall, anchor=west] at (79,33.5) {\textcolor{fslate}{dashed: the probe's boundary}};
\pgfresetboundingbox
\path (0,-2) rectangle (150,103);
\end{tikzpicture}%

\caption{How a probe inherits a dataset-specific feature (schematic).}
\label{fig:spurious}
\end{figure}

\emph{4. Training lies unlike test lies.} Off-policy data suffice for behaviours visible in the text
but ``can harm on-policy performance for `text-ambiguous' behaviours'' such as deception
\cite{kirch2025offpolicy}, and the instructions used to build the training set explained 70.6\% of the
variance in deception-probe performance in one study \cite{natarajan2026targeted}.

\emph{5. Different concepts.} Deception is heterogeneous: probes matched to a type of deception
outperform a universal one \cite{natarajan2026targeted}, and no single direction captures it
\cite{kumar2026pressure}. Domain-adaptation theory bounds a classifier's error on a new domain by its
error on the old one, a divergence between the domains, and the error of the best classifier for both
\cite{bendavid2010theory}; when ``lie'' means different things in two datasets, that last term is large.

\emph{6. Reversed correlates and thresholds.} Falling \emph{below} chance requires a feature whose
correlation with the label reverses between the datasets, as colour does in Colored MNIST, where
empirical risk minimisation reaches 87.4\% training but 17.1\% test accuracy
\cite{arjovsky2019invariant}. Probes invert this way on negated
facts \cite{levinstein2025still}, on expectation-inverted lying \cite{ying2026truthfulness} and on \hpkr{}
\cite{kretschmar2025liars,moustafa2026beyond}. And even when ranking survives, a threshold set on
control data can misfire (\cref{sec:prelim-metrics}).

\emph{Remedies.} Remedies follow the mechanisms: diversify the training data
\cite{liu2024universal,ying2026truthfulness}, treat datasets as separate environments
\cite{arjovsky2019invariant}, or constrain the direction to the difference of means
\cite{marks2023geometry,mallen2024eliciting}, to selected principal components \cite{yoo2026subspace},
or to persona coordinates \cite{mahadik2026persona}. This thesis takes the last route. A direction fitted on
no evaluation dataset cannot absorb a dataset's artefacts (mechanisms 1--3), but it can still miss a kind
of lie its roles do not tell (5), still needs a threshold (6), and is no safer against evasion.

\emph{Learned subspaces.} A subspace learned from labels risks the same failures. \emph{Iterative
nullspace projection} (INLP) \cite{ravfogel2020null} repeatedly trains a linear classifier and projects
its direction out of the data, yielding a basis for the subspace in which a property is linearly
decodable; LEACE removes such information in closed form \cite{belrose2023leace}. Projecting
\emph{onto} an INLP subspace keeps only the directions that source-trained classifiers found, shortcuts
included.

\subsection{Reading what a model knows}
\label{sec:review-knowledge}

A lie, as this thesis defines it, is a statement the model believes to be false
(\cref{sec:background}), so a lie about knowledge is a denial by a model that knows. Several lines of
work show that a model's internal states record whether it knows an answer.
\textcite{kadavath2022language} train language models, with an added value head, to predict the
probability that they can answer a question correctly, which they call P(IK); the prediction is good
within a task and partly transfers to others, although it is poorly calibrated on new tasks, and it
rises when relevant source material is placed in the context. \textcite{slobodkin2023curious} show
that three instruction-tuned models (Flan-T5-XXL, Flan-UL2 and OPT-IML) encode whether a question is
answerable even when they go on to invent an answer: a logistic-regression probe on the first
generated token's last-layer representation separates answerable from unanswerable questions, it
transfers between question datasets, and erasing the subspace it finds with LEACE
\cite{belrose2023leace} lowers the model's own detection of unanswerable questions.
\textcite{gottesman2024estimating} predict how much a model knows about an entity from its internal
representation of the entity alone, before any text is generated. \textcite{ferrando2025know} find,
with sparse autoencoders, directions that detect whether the model recognises an entity, and steering
along them makes a chat model refuse questions about entities it knows, or invent attributes of
entities it would otherwise refuse to answer about.

These results support the premise of the knowledge direction of \cref{sec:k-design}: that whether the
model knows an answer is linearly readable. Our use differs in three ways. The direction is read at the
tokens of a denial written for the model, so the words are fixed and only the knowledge varies. It is
a difference of means over questions marked as known or unknowable, not a classifier trained on the
model's correctness. And it is used to detect lies, not to predict hallucination. The same results also
raise a question that \cref{sec:limitations} returns to: a direction that separates known from
unknowable questions may read whether a question is familiar or answerable rather than whether the
model holds its answer.

\section{Summary of Key Findings}
\label{sec:lit-summary}

\Cref{tab:lit} places the main strands side by side against the properties this thesis needs: a
detector that is cheap to run, needs no labels from the setting in which it is used, and works on the
hardest lies.

\begin{table}[t]
\centering
\small
\caption{How prior approaches meet the requirements of this study. ``Labels'' means labels from
the distribution the detector is evaluated on.}
\label{tab:lit}
\begin{tabularx}{\linewidth}{@{}>{\raggedright\arraybackslash}p{3.3cm}>{\raggedright\arraybackslash}X>{\raggedright\arraybackslash}p{2.3cm}>{\raggedright\arraybackslash}p{2.6cm}@{}}
\toprule
Approach & What it established & Cost per response & Evaluated on \hpkr{} at 70B? \\
\midrule
Truth probes \cite{marks2023geometry,buerger2024truth} & Truth is linearly readable; transfer breaks
on negation and task & one dot product & no \\
Instruction-pair probes \cite{zou2023representation,goldowskydill2025detecting} & Probes read
honesty; transfer is uneven & one dot product & yes: inverted \\
Black-box detectors \cite{pacchiardi2023catch,kretschmar2025liars} & No weight access needed &
one or more model calls & yes: judge fails \\
Upper-Bound Probe \cite{kretschmar2025liars} & Ceiling with in-distribution labels & one dot
product & yes, but uses \hpkr{} labels \\
Truthfulness spectrum \cite{ying2026truthfulness} & General and domain-specific directions coexist &
n/a & no \\
Persona coordinates \cite{mahadik2026persona} & Label-free directions transfer at 3B and 8B & one
dot product & no \\
Self-knowledge readouts \cite{kadavath2022language,slobodkin2023curious,gottesman2024estimating,ferrando2025know}
& Models encode whether they know an answer or can answer a question & one dot product or a small head
& no \\
Transfer analyses \cite{kirch2025offpolicy,natarajan2026targeted,yoo2026subspace} & Transfer
depends on data and on which directions are read & n/a & no \\
\bottomrule
\end{tabularx}
\end{table}

Six findings emerge from the literature.
\begin{enumerate}[leftmargin=1.6em, itemsep=2pt]
\item \emph{Linear readouts are justified, in a weak form.} Many concepts, truth among them, are
  linearly represented, but linearity depends on the variable, and the direction that reads a concept
  is generally not the one that writes it. Differences of means carry all of a binary concept's
  linearly available information.
\item \emph{Truth and deception are linearly readable, but a probe inherits its training
  distribution.} Directions of every generality coexist; a probe trained on one domain mixes them and
  fails where its domain-specific part does not carry over.
\item \emph{The reasons are known.} A probe in thousands of dimensions is underspecified, latches onto
  whatever separates its training set most easily, and inverts where that feature's relation to lying
  reverses.
\item \emph{The hardest case is unsolved without privileged labels.} On \hpkr, no published method
  that does not train on \hpkr{} detects its lies at the benchmark's threshold, and in the published
  evaluations nearly every probe type inverts.
\item \emph{Label-free persona directions are promising but untested where it matters.} They transfer
  better than raw probes on small models but had not been evaluated at 70B or on \hpkr, and a bank
  covers only the kinds of lie its roles tell.
\item \emph{What a model knows is itself readable.} Hidden states record whether the model knows an
  answer or can answer a question, which suggests a label-free signal for lies about knowledge;
  whether it detects them is one of the questions this thesis tests.
\end{enumerate}

The gap this thesis addresses is the intersection: a detector with the cost of a probe and no
evaluation labels, tested on a frontier-size model and on the dataset where existing detectors fail at
the benchmark's threshold. If a probe carries its training domain's directions with it
\cite{ying2026truthfulness} and matching a detector to the type of deception helps
\cite{natarajan2026targeted}, a persona bank with no role that lies as \hpkr's model does should fail
there. That is what \cref{chap:results} finds. The knowledge-denial bank of \cref{chap:method} raises
\hpkr's ranking, but \cref{sec:hpkr-tests} shows that it does so by reading the wording of a denial;
the knowledge direction of \cref{sec:k-design}, built from what the model knows rather than from
role-play, is our answer to that failure.
